\clearpage
\hypertarget{munich-re-transition-disclosure}{%
\section{What an insurer's transition disclosure reveals}\label{method-insurance-munich-transition}}

Munich Re's 2023 report shows how much a public disclosure can resolve.
Its consolidated statements apply IFRS as adopted by the European Union,
together with the stated German commercial-law requirements. The first
application of IFRS~17 is 1 January 2023, with transition at 1 January 2022
and restated insurance comparatives. Financial-instrument comparatives
generally remain under IAS~39, subject to the disclosed classification
overlay. A single year label therefore does not imply a single measurement
basis for every line in the comparative balance sheet
\citep[pp.~155 and 183--184]{munichre2023annual}.\footnote{The retained
English report states that it is an informational translation and the German
original is binding (p.~1). The discussion attributes the disclosed policies
and amounts; the historical AASB texts above are comparison sources.}

The report identifies full retrospection for certain property-casualty
contracts and parts of life reinsurance. Where historical information is
insufficient in life reinsurance, it mostly uses fair-value transition; it
also uses that method for specified property-casualty reinsurance groups and
German primary business. Modified retrospection and fair value together
account for over 90\% of net insurance liabilities at transition. That is a
statement about a liability denominator at a particular date. It supplies
neither a count of contracts nor the assignment and supporting history of
every group \citep[p.~183]{munichre2023annual}.

The opening-equity bridge makes the scale of the different effects visible.
The reported margin component is much larger than the net equity reduction
because the other measurement changes also enter the reconciliation. The
following amounts reproduce the report's signs and billion-euro scale
\citep[p.~184]{munichre2023annual}.

\begin{center}
\begin{tabular}{lr}
\toprule
Opening-equity reconciliation at 1 January 2022 & EUR billions \\
\midrule
Equity under IFRS 4 & 30.9 \\
Modified measurement basis & 25.3 \\
Contractual service margin & $-22.3$ \\
Risk adjustment for non-financial risk & $-5.6$ \\
Loss component & $-1.4$ \\
Deferred taxes & 1.5 \\
Reported change in equity & $-2.5$ \\
Equity under IFRS 17 & 28.5 \\
\bottomrule
\end{tabular}
\end{center}

The five displayed change components sum to $-2.5$ billion, whereas the
displayed endpoints differ by $-2.4$ billion. The report warns that rounding
can cause such deviations; these figures are compatible under rounding to
the displayed precision. The representation should retain both the reported
endpoints and the bridge, with their precision, instead of changing one to
force an exact identity \citep[pp.~1, 155 and 184]{munichre2023annual}.

\clearpage
The separate margin tables also clarify which contracts are being followed.
At 1 January 2022 the reported issued-contract CSM is $-22{,}688$ million
and the held-reinsurance net cost/gain is $368$ million. Adding those reported
totals gives $-22{,}320$ million, consistent with the rounded $-22.3$ billion
margin component above. The issued-contract transition table follows groups
in force at transition, and its CSM amounts concern the general measurement
model and VFA; the associated revenue also includes PAA contracts. Its
later balances therefore cannot be substituted for the whole subsequent
portfolio or used as if CSM and revenue had identical populations
\citep[pp.~210--211 and 254--255]{munichre2023annual}.

Finally, disclosure of a method does not supply all its inputs. Munich Re
reports a 6\% cost-of-capital rate for its non-financial risk adjustment,
applied to projected capital requirements with diversification and discounting.
Its equivalent one-year confidence level is about 90\% at transition, compared
with 88\% for 2023 and 87\% for the comparative period. These are dated
properties of the insurer's calculation, not universal IFRS parameters.
The rate also cannot be imported into the separate fair-value calculation
merely because that discussion mentions cost of capital. Reconstructing an
individual group still requires its cash flows, capital projections, curves,
allocations and supported transition-method choice
\citep[pp.~165, 183 and 185]{munichre2023annual}.
